\chapter{HPV Infection}
\index{HPV Infection}

\section*{Definition and Overview}
Human papillomavirus (HPV) is the most common sexually transmitted infection in the world. It is a group of more than 200 related viruses, of which around 40 infect the genital area. Most HPV infections cause no symptoms and clear on their own within one to two years. However, some types cause genital warts, and around 15 "high-risk" types, particularly HPV 16 and 18, can cause cell changes that lead to cervical, anal, penile, vaginal, vulval, and throat cancers. Vaccination and screening protect against the cancers, and most people clear the infection without ever knowing they had it.

\section*{Epidemiology}
Around 80\% of sexually active people acquire HPV at some point, and at any time around one in ten women has an active genital HPV infection. The highest rates are in people in their late teens and twenties. In many countries, the HPV vaccination program has dramatically reduced the prevalence of vaccine-type HPV and cervical abnormalities in vaccinated cohorts, and some countries are on track to virtually eliminate cervical cancer. Genital warts affect around 1--4\% of sexually active adults and have fallen sharply since vaccination began.

\section*{Aetiology and Pathophysiology}
\begin{itemize}
    \item \textbf{Transmission:} skin-to-skin contact, including vaginal, anal, and oral sex; condoms reduce but do not eliminate transmission
    \item \textbf{Low-risk types:} HPV 6 and 11 cause most genital warts
    \item \textbf{High-risk types:} HPV 16 and 18 cause most HPV-related cancers
    \item \textbf{Pathophysiology:} the virus infects basal epithelial cells; in most people the immune system clears it within months to two years
    \item \textbf{Persistence:} if high-risk HPV persists, it can integrate into host DNA and drive pre-cancerous cell changes
    \item \textbf{Progression:} persistent high-risk infection can progress from cell changes to cancer over many years
    \item \textbf{No treatment for the virus itself:} management targets the warts and cell changes it causes
\end{itemize}

\section*{Clinical Features}
\begin{itemize}
    \item Most infections cause no symptoms and are cleared unnoticed
    \item Genital warts: small, flesh-coloured bumps on the genitals, anus, or surrounding skin
    \item Warts may be flat, raised, single, or clustered
    \item Cell changes are detected by cervical screening, not by symptoms
    \item Advanced changes can cause abnormal bleeding, though this is uncommon
    \item Throat and anal infections are usually silent
    \item Most people never know they have had HPV
\end{itemize}

\section*{Differential Diagnosis}
\begin{itemize}
    \item Genital warts are distinguished from skin tags, molluscum, and other bumps
    \item Cervical cell changes are graded by screening results
    \item Persistent infection is identified by repeat testing, not by symptoms
    \item Other STIs should be considered and tested
    \item Most abnormal screening results resolve without treatment
\end{itemize}

\section*{When to Seek Medical Care}
Women should attend cervical screening every 5 years from age 25 to 74, and anyone with genital bumps, unusual bleeding, or concerns about STIs should be assessed.

\textbf{Contact your local emergency services for:}
\begin{itemize}
    \item Heavy or persistent abnormal vaginal bleeding
    \item Severe pelvic or abdominal pain
    \item Bleeding with collapse or faintness
    \item Any emergency symptoms; HPV itself does not cause acute emergencies
\end{itemize}

\section*{Diagnosis and Investigations}
\begin{itemize}
    \item \textbf{Cervical screening:} the primary test for HPV-related cervical cell changes; samples are tested for HPV, with cell review when positive
    \item \textbf{HPV testing:} detects high-risk HPV DNA
    \item \textbf{Colposcopy:} detailed examination of the cervix after abnormal screening
    \item \textbf{Biopsy:} confirms the grade of cell changes
    \item \textbf{Clinical examination:} for genital warts
    \item \textbf{Other sites:} anal and throat screening for at-risk groups
\end{itemize}

\section*{Management}
\begin{itemize}
    \item \textbf{Observation:} most infections and low-grade changes resolve on their own
    \item \textbf{Wart treatment:} creams, freezing, or laser for genital warts
    \item \textbf{Monitoring:} repeat screening at recommended intervals
    \item \textbf{Colposcopy and treatment:} for higher-grade cervical changes, with loop excision or other procedures
    \item \textbf{Treatment of cancers:} surgery, radiotherapy, and chemotherapy as appropriate
    \item \textbf{Vaccination:} prevents infection with the types in the vaccine, even after other exposures
    \item \textbf{Partner notification:} informing partners is appropriate, though most have already been exposed
\end{itemize}

\section*{Complications}
\begin{itemize}
    \item Genital warts with discomfort and psychological distress
    \item Persistent infection and progression to pre-cancer
    \item Cervical cancer, still a leading cancer in women in countries without screening
    \item Anal, penile, vaginal, vulval, and throat cancers
    \item Recurrence of warts after treatment
    \item Anxiety and stigma around an extremely common infection
\end{itemize}

\section*{Prognosis and Outcomes}
The vast majority of HPV infections clear without consequences, and even persistent cell changes are highly treatable when detected by screening. Cervical screening has reduced cervical cancer deaths by more than half in many countries, and vaccination is on track to make cervical cancer a rare disease. Genital warts are controlled with treatment. The combination of vaccination, screening, and treatment means HPV-related cancer is now largely preventable.

\section*{Prevention and Self-Care}
\begin{itemize}
    \item HPV vaccination for girls and boys, ideally before sexual activity begins
    \item Attend cervical screening every 5 years from age 25 to 74
    \item Use condoms, which reduce transmission
    \item Limit the number of sexual partners
    \item Do not smoke, as smoking increases the risk of persistent infection and cancer
    \item Seek review of any genital bumps or abnormal bleeding
    \item Follow up abnormal screening results promptly
    \item Discuss vaccination for adults where recommended
\end{itemize}

\section*{Sources and Further Reading}
The following reputable sources provide further information for healthcare students and patients seeking reliable, current advice about HPV infection.

\begin{itemize}
    \item Healthdirect. \textit{Human papillomavirus (HPV).} https://www.healthdirect.gov.au/hpv
    \item Cancer Council. \textit{HPV and cancer prevention.} https://www.cancer.org.au/
    \item Australian Government Department of Health. \textit{National HPV vaccination program.} https://www.health.gov.au/
    \item Family Planning NSW. \textit{HPV information.} https://www.fpnsw.org.au/
    \item World Health Organization. \textit{Human papillomavirus fact sheet.} https://www.who.int/
    \item Centers for Disease Control and Prevention. \textit{HPV infection information.} https://www.cdc.gov/
\end{itemize}
